\chapter{PIANTANDO PATATE}\label{piantando-patate}

\hfill\break

Il primo incarico che l'Unef ebbe per me fu, a dispetto di quello che il
generale Meers mi aveva detto, volare qua e là e tenere discorsi. Durò
due settimane e, ragazzi, fu imbarazzante per tutte le persone
coinvolte. Per quanto mi riguarda, era imbarazzante indossare la mia
nuova uniforme davanti a una folla di soldati, persone che erano state
in combattimento come me, in situazioni peggiori di quelle in cui mi ero
trovato io e indossavano ancora le stesse uniformi e svolgevano ancora
gli stessi compiti, mentre io indossavo aquile d'argento, volavo in giro
a bordo di una Poiana lucente e mangiavo buon cibo con gli alti
ufficiali. Mi faceva sentire un ciarlatano. Per i soldati, in attesa che
io facessi il mio discorso preconfezionato, era imbarazzante perché
tutti sapevano che non meritavo di essere un ufficiale, non me l'ero
guadagnato, ma non potevano obiettare nulla al riguardo. Fu uno schifo.
Quando il mio discorso grazie a Dio finiva, facevo domande alle persone,
lasciavo che parlassero delle loro esperienze e ignoravo l'imbarazzo del
mio essere lì. Funzionava per tutti. Quando stavo in mezzo alla folla,
cazzeggiando coi soldati, lasciando che fossero loro a raccontare le
storie di quello che avevano fatto durante la fallita invasione dei
ruhar, potevo essere soltanto Joe Bishop, e loro potevano essere
soldati, e potevamo parlare.

E poi dovevo tornare sulla Poiana, volare alla base successiva e rifare
tutto da capo. Lo odiavo.

Perciò fu una benedizione quando l'Unef mi trovò un vero compito:
piantare patate. I kristang avevano detto all'Unef che era ora che gli
umani si coltivassero da sé il cibo su Paradiso, per ridurre il loro
carico logistico e darci un po' di sicurezza alimentare, nel caso in cui
le azioni della flotta di sopra interrompessero le spedizioni. Qualche
genio del quartier generale dell'Unef doveva aver deciso che i miei
discorsi avevano fatto il loro tempo, e, letto il mio fascicolo
personale e visto che ero del Nord del Maine, aveva avuto la brillante
intuizione che dovevo essere un esperto nella semina delle patate. Il
mio nuovo compito consisteva nel coordinare le attività agricole in un
quarto del continente occupato dall'Unef. Non piantavamo soltanto
patate, ovvio, avevamo in programma di crescere un'ampia varietà di
colture da semi spediti dalla Terra. Piantare patate era la dicitura
ironica che usavamo noi, anche se si aspettavano che ne fossi
ufficialmente entusiasta. Piantare patate mi permetteva di volare per
Paradiso, mostrarmi alle truppe e, stando al personale delle pubbliche
relazioni dell'Unef, mi dava l'opportunità di essere visto come un
soldato semplice laborioso, che faceva il suo dovere attraverso il duro
lavoro e l'iniziativa. Era un sollievo fare qualcosa di utile.

Oltre a un corso intensivo in agricoltura, ottenni qualche altro
vantaggio. In quanto colonnello, ora avevo il mio Crivee personale e un
autista, un certo soldato Randall, nella mia nuova base, al confine tra
il settore americano e quello cinese. «Porca troia.» Non potevo crederci
quando Randall mi mostrò il mio mezzo di trasporto. Di fatto, quello che
vidi mi lasciò di stucco. Il Crivee era un veicolo ruhar standard, con
il simbolo Unef su entrambi i lati e sul tetto. E un Barney di peluche
viola alto un metro e mezzo legato alla griglia frontale. «Siamo a mille
anni luce dalla Terra, come cazzo avete fatto voi idioti a procurarvi un
Barney?» Ogni singolo pezzo di equipaggiamento di cui avevamo bisogno
doveva fare il lungo viaggio fino in Ecuador, poi sull'ascensore
spaziale, su una nave kristang, su una nave madre thuranin, attraverso
wormhole, quindi un percorso a ritroso verso la superficie di Paradiso.
Tutta la nostra attrezzatura veniva ispezionata in modo meticoloso, per
sincerarsi di fare il miglior uso possibile di ogni chilo di massa e
metro quadrato di spazio. Eppure, non si sa come, qualche burlone era
riuscito a infilare di soppiatto un gigante Barney imbottito in un
container. Tutto quello che volevo sapere era, per l'amor di Dio,
perché? E cos'altro aveva intrufolato la gente a bordo, per poi
spargerlo in modo clandestino sulla superficie di Paradiso?

«L'abbiamo acquisito tatticamente, signore», disse Randall con una
faccia seria.

«Nel senso che l'avete rubato.»

«A porta aperta anche il giusto vi pecca. Volevamo farla sentire il
benvenuto, signore. Lo scherzo è finito, lo toglierò dalla griglia.»

«No, che cazzo, lascia stare.» Il mio Crivee non era uno scherzo più di
quanto non lo fosse l'idea di me colonnello a tutti gli effetti.
Camminai verso la parte anteriore e ispezionai il dinosauro viola
ghignante. Forse ai bambini criceti sarebbe piaciuto.

\hfill\break

Fu mentre ero fuori a piantare patate che trovammo il nostro primo
biscotto della fortuna. Non avevamo ricevuto comunicazioni dalla Terra
da quando l'ultimo gruppo di umani era arrivato su Campo Alfa e io ero
diretto lì nel terzultimo convoglio. Da quando eravamo atterrati su
Paradiso, non c'era stata alcuna comunicazione né da né verso la Terra;
i kristang ci avevano detto che la situazione militare nello spazio non
permetteva loro il lusso di rimandare sulla Terra umani, neppure feriti
con voli di evacuazione medica. Saremmo rimasti bloccati su Paradiso per
tutta la durata della missione, la cui indeterminatezza metteva a
disagio qualunque umano. Il comando dell'Unef era agitato dal fatto di
non essere in grado di inviare resoconti della situazione sulla Terra, o
di non riceverne ordini e indicazioni. I soldati comuni come me si
preoccupavano della famiglia e degli amici. I governi avevano fatto
progressi nel ripristino dell'elettricità e delle altre infrastrutture?
Rifornimenti e rinforzi erano in arrivo? I ruhar avevano attaccato di
nuovo la Terra? La Bürgermeister mi aveva detto che era molto
improbabile che i ruhar avrebbero organizzato, o potuto organizzare,
un'altra spedizione sulla Terra nei prossimi anni. Sempre che mi stesse
dicendo la verità.

Poi trovammo i biscotti della fortuna e tutto cambiò. I kristang
ispezionavano tutti i rifornimenti consegnati agli ascensori spaziali
sulla Terra, in cerca di merce di contrabbando. Scoprimmo, e ci
inquietò, che i kristang cercavano con particolare attenzione
dispositivi di archiviazione di dati digitali che bruciavano con impulsi
magnetici o ultravioletti. Quello che non si aspettavano era che alcune
persone stampassero parole su carta e incollassero quella carta
all'interno dei pacchi alimentari. Per gli scanner kristang, un
contenitore di cartone rinforzato, con minuscole scritte stampate
all'interno, somigliava a qualsiasi altro contenitore di cartone. Quei
contenitori venivano imballati sulla Terra, consegnati su Paradiso dai
kristang e disfatti dagli umani. Quando i ragazzi della prima fornitura
aprirono un pacchetto di cibo con un biscotto della fortuna all'interno,
grazie a Dio, ebbero l'intelligenza di non dire nulla al riguardo sui
loro zPhone; dopo di che fu tutto un passaparola. I biscotti della
fortuna furono rimossi con cura dalla confezione e consegnati a mano al
quartier generale dell'Unef, dove causarono una tempesta di panico
totale.

Venni a sapere della faccenda dall'agente dell'intelligence della III
divisione, le cui mani stavano di fatto tremando mentre mi spiegava il
tutto. Le notizie provenienti dalla Terra erano cattive, e i biscotti
della fortuna contenevano codici segreti per autenticare i dati al
quartier generale dell'Unef, quindi avevamo la certezza che erano vere.
Dopo che la nostra Expeditionary Force aveva lasciato la Terra, le cose
avevano cominciato a peggiorare. Anche prima di andarmene, avevo sentito
dire che i kristang stavano usando il pugno di ferro: dirigendosi dove
avrebbero dovuto concentrarsi i nostri sforzi di riparazione delle
infrastrutture, impossessandosi di complessi minerari e di raffinazione
per materiali critici e cercando di controllare in modo rigoroso
determinate informazioni su internet. Come la maggior parte delle
persone, avevo pensato che i kristang stessero facendo quello che
dovevano per prepararci nel caso in cui i ruhar fossero tornati, e che
ci sarebbero sempre stati problemi quando due specie aliene avessero
dovuto adattarsi l'una all'altra. E che i civili tendevano a lamentarsi
di tutto in ogni caso.

I biscotti della fortuna ci dicevano che le condizioni sulla Terra erano
peggiorate, molto peggiorate. I kristang si erano impossessati di alcuni
dei migliori terreni agricoli del pianeta, tra cui vaste aree del
Midwest americano, per coltivare i loro raccolti e allevare i loro
animali per il sostentamento. Gli agricoltori non erano stati risarciti
dai kristang per la terra che avevano perso e, nei casi in cui gruppi
avevano cercato di impedire alle lucertole di occupare i loro terreni,
erano stati uccisi a colpi di maser dall'orbita. Ma il peggior esproprio
non era stato quello dei terreni agricoli: i kristang volevano che il
lago Superiore e il mar Caspio producessero i tipi di pesce che
mangiavano loro. Progettavano di sterilizzare quegli enormi laghi con
potenti raggi gamma e di ricreare nell'acqua la propria biosfera. La
gente era stata avvisata del fatto che avrebbe dovuto mantenersi a cento
chilometri di distanza dalle rive di questi specchi d'acqua quando i
raggi gamma avrebbero colpito, cosa che sarebbe accaduta non appena le
imponenti dighe sarebbero state completate e il satellite a raggi gamma
sarebbe stato pronto, nel giro di un anno o giù di lì. Le proteste in
città come Shanghai, Chicago, San Francisco e Parigi erano state
soppresse con attacchi maser, nel momento in cui i kristang avevano
saputo che gli sforzi compiuti dai governi umani per arginare il
dissenso si erano rivelati inefficaci contro ``elementi sovversivi e
traditori''.

I ruhar ci avevano attaccato e noi avevamo pensato che i kristang
fossero i nostri salvatori, ma il salvataggio si era trasformato in
un'invasione. Tutto cambiò e, in un certo senso, nulla cambiò per noi su
Paradiso. Cosa poteva fare l'Unef? Tutti i nostri rifornimenti ci
venivano consegnati dai kristang. Non avevamo modo di tornare a casa
senza i kristang. Anche se ora li consideravamo dei nemici, come
facevamo a sapere che i ruhar erano migliori? I criceti su Paradiso non
erano comunque in grado di aiutarci, quindi non contava che i kristang
ci piacessero o meno.

Gli umani su Paradiso erano fottuti in ogni caso.

\hfill\break

Il colonnello Wilson venne in ufficio mentre stavo finendo una tazza di
caffè e leggendo i rapporti sulle forniture, preparandomi per un'altra
incursione di semina delle patate. Il caffè era amaro e insipido allo
stesso tempo e si era raffreddato, ma ero determinato ad assaporarne
ogni singola goccia. Quei chicchi di caffè avevano viaggiato attraverso
anni luce per entrare nella mia tazza lì su Paradiso. Aveva il sapore di
casa. Non sapevo se ne avremmo mai ricevuto un'altra partita dalla
Terra. Wilson si versò ciò che ne rimaneva nel recipiente in una tazza,
sorseggiò e fece una smorfia: «Sono appena arrivato dal comando del
generale Meers, non sa come dovremmo gestire queste voci sui biscotti
della fortuna. I kristang lo scopriranno di sicuro, mi sorprende che non
se ne siano già lagnati. A meno che non l'abbiano fatto e il quartier
generale lo stia passando sotto silenzio».

Con il loro controllo quasi totale delle nostre comunicazioni, non era
possibile che i kristang non sapessero cosa stava succedendo. Forse non
se ne curavano. La gente sulla Terra sapeva cosa stava succedendo lì, e
che cosa avrebbero potuto farci, con le astronavi kristang in orbita, in
grado di bombardare qualsiasi punto del pianeta con raggi maser, missili
e cannoni a rotaia? «Abbiamo bisogno di cambiare mentalità sulla nostra
missione qui», suggerii. ``Mentalità'' era una parola gettonatissima su
molte delle slide di PowerPoint dell'esercito. «Noi, voglio dire, almeno
noi americani abbiamo pensato che l'attacco dei ruhar fosse un'altra
Pearl Harbor, e di trovarci nel Pacifico del Sud durante la seconda
guerra mondiale, a contrattaccare. Ora sappiamo che è una stronzata. La
verità è che siamo nell'operazione Torcia e non siamo né gli Alleati né
l'Asse, siamo i berberi.»

«Non ti seguo, Bishop.» Da ufficiale dell'esercito degli Stati Uniti,
Wilson doveva sapere che ``operazione Torcia'' era l'espressione in
codice per l'attacco condotto da americani e inglesi contro il Nord
Africa nel novembre del 1942. «Berberi?»

«I berberi.» Stavo per aggiungere ``signore'', questa cosa del grado di
colonnello era ancora nuova per me. «Penso che sia così che li chiamano.
I nativi del Nord Africa, che si trovarono intrappolati tra l'invasione
alleata da un lato e i tedeschi, gli italiani e i francesi di Vichy
dall'altro. A nessuna delle due parti importava un cazzo dei nativi, o
delle loro terre. I tedeschi e gli italiani volevano il Nord Africa per
poter controllare il Mediterraneo e costringere gli inglesi a lasciare
l'Egitto, in modo da avere accesso ai giacimenti petroliferi. Noi e gli
inglesi volevamo respingere Rommel attraverso il Mediterraneo, in modo
da poter colpire la Sicilia e poi l'Italia.» Stavo surfando sulla
verità: gli Alleati non avevano deciso immediatamente dove andare dopo
la vittoria in Nord Africa. «I berberi sono rimasti intrappolati nel
mezzo e tutto quello che hanno potuto fare è stato seguire gli ordini di
qualsiasi parte occupasse il loro territorio in quel momento, togliersi
di mezzo e cercare di rimanere in vita. Ora siamo noi. Non contiamo un
cazzo per nessuna delle due parti, se non per come possono usarci, o
usare la Terra in quanto base di sosta. Siamo solo barrette croccanti
sotto i cingoli dei loro carri armati. Dobbiamo smettere di pensare che
siamo alleati dei kristang e capire che siamo solo truppe native che i
kristang possono usare come carne da cannone. I francesi liberi allora
avevano truppe berbere chiamate Goumier, ma non erano trattati come i
soldati francesi. L'Unef deve smettere di pensare a grandi ambizioni e
concentrarsi sulla sopravvivenza.»

Wilson si accigliò, ma annuì: «Può darsi che tu abbia ragione, Bishop».

Ingoiai la feccia del mio caffè e presi casco e occhiali: «Se vogliamo
avere qualche possibilità di sopravvivenza, è meglio che torni a
piantare patate. Una caterva di patate.».

\hfill\break

Presi una Poiana diretta in un villaggio nel bel mezzo del nulla, ore di
ascolto del ronzio dei motori e terreni agricoli che scorrevano sotto di
noi. Purtroppo, ero rimasto solo con i miei pensieri, e non era un posto
piacevole in cui stare. Quando avevo visto la prima nave d'assalto ruhar
scivolare attraverso il campo di patate nella mia città, la prima cosa
che mi era venuta in mente, subito dopo ``Merda, sta succedendo
davvero'' e ``Perché cazzo gli alieni avrebbero invaso Thompson
Corners'', era stata ``Game over''. A parte le stronzate dei film di
Hollywood, qualsiasi specie dotata della tecnologia, delle risorse e
dell'incentivo necessari per viaggiare tra le stelle avrebbe schiacciato
l'umanità come un insetto. Dimenticatevi le fantasie di impavidi umani
che sconfiggono gli alieni a terra. Questi ultimi potrebbero stazionare
comodamente in orbita e polverizzarci a loro piacimento. Tutta la
determinazione e lo spirito umano nell'universo non servirebbero a un
cazzo se potessero colpirci e noi non potessimo raggiungerli. Anche se
riuscissimo a reindirizzare e lanciare missili Icbm a punta nucleare,
quei missili non potrebbero raggiungere l'alta orbita e gli alieni
dovrebbero essere del tutto ciechi per non vedere il gas di scarico del
razzo. Sarebbe sorprendente se una testata Icbm arrivasse a meno di
centosessanta chilometri da una nave aliena; è probabile che non
riuscirebbe a superare neanche lo strato inferiore dell'atmosfera. Game
over.

Non ero l'unico a denigrare le prospettive dell'umanità di sopravvivere
a un'invasione aliena, è quello che ci avevano detto i kristang. Le
lucertole avevano raccontato un sacco di bugie, ma in questo caso erano
state sincere. Avevano detto la verità, per farci capire quanto avessimo
bisogno di loro per proteggerci dai ruhar, o almeno così pensavamo
all'epoca. Il punto era: umani bloccati a terra, alieni con la
postazione strategica dell'orbita e oltre a disposizione. Game over.

Quando avevo deciso di provare a catturare un soldato alieno, l'unica
cosa cui avevo pensato era che l'umanità avesse bisogno d'informazioni
sugli invasori, se avessimo avuto la possibilità di sopravvivere in
qualche modo. E che dovevo fare qualcosa, anche se era avventato e
stupido. E che, con tutta probabilità, sarei stato ucciso, ma, dato che
l'intera umanità stava per essere spazzata via, Whisky Tango Foxtrot,
no?

Poi, da oltre l'orizzonte, era arrivata la cavalleria kristang ed
eravamo salvi. Cazzo, mi ricordo di quella sensazione, quando quel carro
attrezzi ruhar ci braccava. Mi pisciavo sotto in quel momento, non mi
vergogno ad ammetterlo, perché mi pisciavo sotto e avevo la bocca secca
e mi tremavano le mani, ma resistevo comunque. Mi stavo preparando a
morire in un lampo di fuoco, e poi il cielo aveva brillato di nuovo ed
era avvenuto un miracolo, quando quella navicella ruhar aveva cabrato ed
era sfrecciata in cielo con un boom supersonico.

Eravamo salvi! Tutte le mie paure erano state spazzate via. Ed ero grato
ai nostri soccorritori.

Erano tutte stronzate.

Questo era anche peggio.

\hfill\break

La Poiana atterrò, entrai in un Crivee, che non era il mio Crivee Barney
personale ma un modello generico, e partimmo rombando subito dopo: un
convoglio di piantatori di patate, composto da quattro Crivee blindati
montanari e sei pesanti camion con rimorchio. Mi irritava il fatto che
il convoglio avesse dovuto aspettarmi; odiavo aspettare gli ufficiali
quando ero un soldato semplice, ora odiavo farmi aspettare dai soldati.

Il mio Crivee odorava di piedi, piedi puzzolenti. E qualcosa come
dopobarba o acqua di colonia? Doveva essere stato parcheggiato in un
posto in cui i criceti usavano quel disgustoso fertilizzante puzzolente,
poi qualcuno aveva cercato di deodorarlo con acqua di colonia a buon
mercato. «Cazzo, puzza come l'acqua di colonia di mio nonno», brontolai.

«Mi dispiace, signore», disse la conducente, era un soldato semplice con
scritto ``Park'' sulla targhetta del nome. «Puzzava già così quando
l'abbiamo preso al parco macchine.»

«Puzzava di più questa mattina», disse un soldato di nome Olafson dal
posto del passeggero. «Abbiamo lasciato aperti i finestrini per far
uscire l'odore.»

«Olafson, uhm?» Era biondo, un metro e novantacinque circa. «Sei sempre
stato così alto oppure l'esercito ti ha dato troppo da mangiare?»

«Mia madre era una brava cuoca, signore.»

«Sì, sì.» Non ero di buona compagnia, non ero dell'umore giusto per
essere di buona compagnia, perciò affondai il naso nei rapporti
scaricati sul mio tablet e i due soldati sui sedili anteriori capirono
l'antifona. L'ufficio di intelligence agricola del quartier generale
dell'Unef aveva preparato una serie di rapporti per me sulla prossima
area in cui avremmo piantato colture: test del suolo, clima,
precipitazioni, acque sotterranee, i tipi di colture che i ruhar stavano
facendo crescere lì e le sostanze chimiche che avevano usato. Il fatto
che l'Unef avesse istituito in fretta e furia un ufficio di intelligence
agricola la diceva lunga su quanto fosse scomoda la nostra situazione su
Paradiso. Avevo pensato fosse uno scherzo quando mi avevano affidato
l'incarico di piantare le colture di cui avevamo bisogno per
sopravvivere; avevo aiutato i miei genitori con la loro piccola
fattoria, ma non ero per niente un esperto di agricoltura. Era
sorprendente per me quanto avessi imparato in breve tempo. La paura di
morire di fame è un grande incentivo. Il luogo in cui avremmo piantato
le colture sembrava essere perfetto: era abbastanza vicino all'equatore
perché potessimo ottenere due o tre raccolti all'anno, c'erano un sacco
di precipitazioni e buone acque sotterranee per l'irrigazione durante la
stagione estiva secca. Controllai i manifesti di carico dei camion nel
nostro convoglio, eravamo così stracarichi di fertilizzante perché la
terra locale doveva essere preparata prima che gli organismi terrestri
potessero crescervi bene. Quando avevo letto tutto quello che potevo
capire sull'agricoltura, mi ero dedicato agli altri rapporti che l'Unef
si aspettava che leggessi. Chi l'avrebbe immaginato che essere un
ufficiale avrebbe comportato così tanta lettura? Era come ritornare a
scuola. Speravo che non ci fosse un quiz il giorno successivo.

Viaggiammo in silenzio per un po', io esaminando i rapporti, Park e
Olafson bisbigliandosi qualcosa di tanto in tanto sui sedili anteriori.
La nostra posizione di terzo veicolo nel convoglio faceva sì che ora
dovessimo chiudere i finestrini per tenere fuori la polvere, ora aprirli
per far uscire il calore. La mattina era stata fredda, adesso la
temperatura si stava riscaldando in fretta, mentre attraversavamo il
terreno agricolo pianeggiante. La nostra destinazione era un'area
evacuata di recente dai ruhar, dove avevamo in programma di rimuovere
dal terreno i resti dei loro raccolti abbandonati, preparare il
terriccio e piantare i nostri semi. Evacuata di recente ed evacuata a
denti stretti. La zona era stata sotto la responsabilità dell'esercito
indiano; secondo i rapporti che stavo leggendo, avevano dovuto rimuovere
con la forza le famiglie di criceti dalle loro case e c'era stato uno
scontro a fuoco in cui un soldato indiano era morto e due erano rimasti
feriti in modo grave. L'intera situazione su Paradiso era precaria; i
kristang mantenevano ancora la posizione strategica, anche se con un
solo cacciatorpediniere e una fregata in volo. I ruhar venivano ancora
evacuati fuori dai confini del pianeta, spesso opponevano resistenza,
con crescenti incidenti di sabotaggio contro l'Unef. C'era affollamento
alla stazione base dell'ascensore, l'Unef vi aveva dovuto allestire un
campo profughi, e le operazioni logistiche necessarie a mantenere i
criceti riforniti e sotto controllo stavano distraendo il quartier
generale. Le navi da trasporto ruhar continuavano a raggiungere la cima
dell'ascensore per portare via i criceti, arrivavano in ritardo, senza
più un programma regolare e in numero minore, scortate da una piccola
protezione di navi da guerra kristang e caricate in modo frenetico e
precipitoso dalle lucertole, che sospettavano un'altra incursione ruhar.
A terra, l'Unef si preoccupava più della sopravvivenza, in particolare
del nostro approvvigionamento alimentare, che di cacciare fino
all'ultimo criceto fuori dai confini del pianeta. Le comunità di criceti
erano passate dalla cooperazione alla riluttanza, alla resistenza attiva
e al sabotaggio. Il sabotaggio si spinse oltre la manomissione dei
camion che avrebbero dovuto essere utilizzati per evacuare i ruhar dai
loro villaggi; i criceti fecero saltare i ponti, rendendo difficile o
impossibile l'evacuazione via terra. L'Unef iniziò a utilizzare chiatte,
pensando che nessuno avrebbe potuto far esplodere un fiume, ma il
trasporto fluviale era lento e prendeva una via traversa, e un viaggio
più lungo richiedeva una quantità maggiore di rifornimenti. Far
viaggiare in volo i criceti era un'opzione costosa e, con l'aumento dei
voli, gli aerei avrebbero avuto bisogno di più tempo per la
manutenzione, creando un ciclo discendente della disponibilità degli
aeromobili. I rapporti che avevo letto non fornivano buone risposte su
quale fosse il modo migliore per continuare la missione di evacuazione,
l'unica cosa certa era che l'Unef aveva bisogno di tenere i criceti in
movimento costante, il più rapido possibile. Più terra veniva liberata
dai criceti, più l'Unef poteva concentrare gli sforzi ed esercitare un
controllo più stretto sulla popolazione ruhar restante.

Agenda e disponibilità della manutenzione delle attrezzature. Era
qualcosa di cui non mi ero mai dovuto preoccupare prima. Ora che ero un
ufficiale, un alto ufficiale, avrei dovuto conoscere queste stronzate.
Conoscerle, considerarle, pensarci, trovare soluzioni e implementarle.
Io.

Non ero soltanto una trovata pubblicitaria, ero un colonnello
dell'esercito americano a tutti gli effetti. Con l'autorità e la
responsabilità di un colonnello a tutti gli effetti. Avevo bisogno di
essere un colonnello a tutti gli effetti.

«Park, Olafson, siete di stanza qui da molto tempo?»

«Da quando siamo arrivati qui, signore», disse Olafson.

«Fatemi un resoconto della situazione.»

«Un resoconto della situazione, signore?», chiese Park dal posto di
guida. Potevo vedere i suoi occhi nel retrovisore.

Sapeva cosa fosse un resoconto della situazione, così esplicitai la mia
richiesta: «Che cosa è successo con i criceti da queste parti?».

Olafson si girò in modo goffo sul sedile, così poteva guardarmi mentre
parlavamo. «Non molto, signore. Roba di basso livello, qualche
sabotaggio che ha rallentato la tabella di marcia del programma di
evacuazione, ma nessun vero problema, nessuna violenza. Andavamo molto
d'accordo con i nativi, finché, lei, ehm», lanciò un'occhiata a Park,
«ha buttato le loro navi giù dal cielo, signore», disse quest'ultima
cosa con uno sguardo di ammirazione così evidente che ne provai
imbarazzo.

«Io non ho fatto niente, Olafson. C'era una squadra di soldati a
difendere il Vettore. Non era l'unico punto nevralgico nel pianeta quel
giorno.»

«Ehm, no, signore. Era piuttosto tranquillo qui quel giorno, non c'è
nulla di molto strategico per cui valga la pena lottare da queste parti.
Se non avessimo sentito parlare dell'attacco sui nostri zPhone e visto
l'azione in corso in cielo, sarebbe stato un giorno normale per noi.»

Park annuì dal sedile del conducente: «Tutto liscio, signore. Quel
giorno abbiamo indossato l'armamentario di battaglia e ci siamo
preparati a resistere, ma non è successo niente, poi ci hanno dato il
cessato allarme. Questa parte del pianeta è in mezzo al nulla, davvero
in culo a Nettuno. Ora, più avanti, dove pianteremo le patate o
qualunque altra cosa, i criceti hanno dato un sacco di problemi
all'esercito indiano. Nessuna sparatoria qui, nessun morto da entrambe
le parti, ma quei cazzo di criceti\ldots» Potevo vedere la sua
esitazione nel retrovisore.

«Va tutto bene, Park, sei un soldato, cazzo, puoi parlare come tale.»

«Sissignore. I criceti hanno opposto resistenza in ogni modo possibile.
L'equipaggiamento era stato sabotato, perciò gli indiani non potevano
usare nulla qui, tutta la loro attrezzatura è stata portata da fuori o
ricostruita. I rifornimenti d'acqua erano contaminati, i criceti si
muovevano in modo da incasinare il programma di evacuazione. Le cose si
sono complicate così tanto che gli indiani hanno dovuto mettere l'intera
regione in isolamento, tipo agli arresti domiciliari, e quando gli
autobus si fermavano per portare via i criceti, quelli non lasciavano le
loro case. Si sedevano sul pavimento e dovevano essere trasportati fuori
a forza e caricati sugli autobus. Rallentavano tutto rispetto al
programma. Le luc\ldots{} i kristang, volevano vaporizzare un villaggio
dall'orbita, in modo che i criceti ricevessero il messaggio che non
potevano fare i furbi con noi oltre un certo limite. L'Unef pensava che,
se i kristang fossero stati costretti a intervenire, avrebbero visto la
nostra missione qui come un fallimento, perciò gli indiani hanno
schierato un battaglione britannico come rinforzo\ldots»

«Tac! Questo è stato un bel calcio in culo, eh, signore?», sogghignò
Olafson. «Gli inglesi sotto il comando degli indiani?»

Annuii in silenzio. L'Unef sottolineava l'importanza della cooperazione
internazionale, e gli ufficiali avevano ricevuto severe istruzioni di
non incoraggiare o consentire rivalità nazionaliste, ma tali ordini
avevano limiti pratici. «Serve da esempio, siamo tutti umani quaggiù»,
dissi, in un patetico tentativo di seguire lo spirito delle regole.

«Sì, sì», disse Park, «tra di loro hanno sistemato il posto in tempo, ma
la cosa è stata deleteria per la tabella di marcia nel settore
britannico, perciò ora si stanno mettendo in pari.» Il programma
accelerava perché ci avvicinavamo al giorno dell'evacuazione, quello in
cui l'ultimo criceto sarebbe salito sull'ascensore spaziale. A mano a
mano che sempre più luoghi di Paradiso venivano sgombrati dai ruhar,
potevamo concentrare le nostre forze e far correre i criceti. Il
programma all'inizio sembrava illogico, perché dava la priorità
all'evacuazione di alcune aree poco popolate del pianeta, e l'Unef aveva
ipotizzato che i kristang volessero che quelle aree fossero ripulite per
prime in modo da poter piantare le loro colture. Ma le lucertole non
fecero una mossa per preparare i terreni agricoli abbandonati e
mandarono solo piccole squadre a ispezionare le aree. Girava voce che
secondo le lucertole quelle aree contenessero reliquie della
super-civiltà antica degli Anziani, reliquie che erano il vero premio
per cui valesse la pena di combattere per il controllo di Paradiso. Del
resto, l'intero pianeta era occupato da terreno agricolo generico, che
abbondava anche nel Braccio di Orione della galassia. Forse. Tutto
quello che importava all'Unef era che i kristang ci avevano detto di
liberare prima quelle aree, così lo facemmo.

«Stiamo arrivando a Habitrail, signore, dritto davanti a noi», indicò
Olafson. La strada saliva sulla cresta di un'altura, davanti potevamo
vedere un'ampia valle fluviale poco profonda, la strada che passava
sopra un ponticello davanti a noi e poi tra le case sparpagliate di una
città di una certa dimensione. Tipici appezzamenti di terreno agricolo,
un po' di terra nuda fresca di raccolto, un po' di terra dal colore
giallo dorato dell'equivalente del grano che i criceti piantavano, e
imponenti silos di cereali. Il sottile nastro luminoso dei binari della
ferrovia andava dritto come una freccia da sud-est a nord-ovest,
formando un angolo col nostro percorso, e, in lontananza, erano
parcheggiate una locomotiva elettrica e una dozzina di vagoni merci. Il
piccolo treno era bloccato lì; i sabotatori avevano distrutto o minato i
binari e fatto saltare o indebolito i ponti ferroviari nella regione. Il
piano originale dell'Unef contava sulle ferrovie per trasportare in
fretta un gran numero di criceti. Il problema era che questi l'avevano
capito subito e ogni linea ferroviaria del pianeta era stata resa
inservibile in un paio di settimane.

«Habitrail?» Risi.

Olafson fece un ampio sorriso. Non era più così intimorito o intimidito
da me: «Il nome criceto suona come ``Hah-bah-tahlin'', quindi quei
burloni di prima categoria hanno messo il cartello ``Habitrail'\,''. Il
sindaco criceto qui si è incazzato quando ha scoperto che sulla Terra
Habitrail è il marchio di una gabbia per criceti». Olafson non era
empatico. «Che si fotta.»

\hfill\break

Il mio zPhone emise un suono attutito, qualcosa come: «Planter, qui è
Stinger guida».

«Qui Planter», risposi, dopo aver estratto il mio zPhone dal mucchio di
giubbotti antiproiettile sotto cui l'avevo sepolto. Planter era il mio
attuale indicativo di chiamata, mi era stato assegnato da qualche
burlone nel quartier generale dell'Unef e il mio assistente part time si
era incazzato e aveva cercato di farlo cambiare in qualcosa che suonasse
più figo, ma a me piaceva. ``Piantatore'', se i criceti stavano
ascoltando le nostre comunicazioni, e dovevo presumere che potessero,
era un bell'indicativo di chiamata, amichevole, non aggressivo, che si
sperava enfatizzasse la natura pacifica delle missioni che mi erano
state assegnate. Piantavamo semi solo sulla terra da cui i criceti erano
già stati evacuati, le nostre missioni non sottraevano loro alcuna
risorsa, e noi evitavamo il contatto con le popolazioni ruhar ogni volta
che potevamo. A differenza di quel giorno, perché il nostro convoglio
doveva attraversare Habitrail, che si trovava sull'unica strada che
portava alla nostra destinazione. «Vai avanti, Stinger capo.» Stinger,
``Pungiglione'', diceva il mio pacchetto informativo per la missione,
era l'indicativo di chiamata della scorta di due navi che avrebbero
fornito supporto aereo mentre attraversavamo Habitrail. Un paio di
cannoniere Pollo armate fino ai denti avrebbero dovuto dissuadere
qualsiasi criceto dal fare l'intrepido mentre passavamo. Io, se vedessi
un paio delle nostre stesse cannoniere usate contro di noi, sarei più
arrabbiato che intimidito. Abbastanza arrabbiato da agire. Se la rivalsa
era certa? Con tutta probabilità no. In ogni caso, era bello avere un
paio di cannoniere a sostenerci.

«Planter, sappi che il mio braccio destro ha un calo di potenza in un
motore e deve tornare alla base.»

Ora che la voce non proveniva da sotto uno strato di giubbotti
antiproiettile, potevo riconoscere che il pilota era una donna, con un
accento americano del Midwest. Mi aspettavo che dicesse ``Puoi
contarci!''.

«Ricevuto, tornando alla base, Stinger, puoi coprirci tu attraverso il
posto di blocco Mike 2?» M2 era la designazione di Habitrail sulla
mappa.

«Affermativo, Planter, vi coprirò da ponte a ponte, passo.» Intendeva
dal ponte dal nostro lato di Mike 2 al ponte dall'altra parte della
città. Oltre il secondo c'erano terreni agricoli aperti, i potenziali
punti di pericolo in cui avremmo potuto aver bisogno di copertura aerea
erano i passaggi obbligati dei due ponti e la città stessa.

«Affermativo, Planter. Farò un volo radente sulla città per svegliarli.»

«Ricevuto, Stinger. Planter chiudo.» L'Unef aveva discusso se fosse
meglio che i convogli arrivassero nelle città senza preavviso, o che la
copertura aerea facesse prima un sorvolo. Il dibattito sull'argomento
era stato breve, perché un convoglio che viaggiava su strade sterrate,
attraverso terreni agricoli pianeggianti, difficilmente poteva
aspettarsi di avvicinarsi di soppiatto a qualcosa: il pennacchio di
polvere sulla scia di un convoglio era un pugno nell'occhio, visibile a
chilometri di distanza.

Il primo ponte che superammo era un brutto blocco di cemento, costruito
dai kristang la prima volta che avevano occupato il pianeta. Due soldati
della squadra di sicurezza locale ci salutarono con la mano, erano di
stanza sul ponte per assicurarsi che nessuno l'avesse manomesso
dall'ultima ispezione. Questa era la situazione che affrontavamo; se
c'era una cosa che i criceti potevano sabotare, l'avrebbero fatto.
Bastava distogliere lo sguardo un paio d'ore perché un ponte venisse
fatto esplodere o danneggiato in qualche modo per renderlo inservibile,
una chiatta fluviale venisse affondata, ripetitori venissero abbattuti,
motori di veicoli incustoditi venissero bruciati. Con mia grande
sorpresa, attraversammo la città senza incidenti, nemmeno un piccolo
criceto dalla testa calda che ci lanciasse una pietra, o una zolla di
fango, o un pezzo di letame. Di solito potevamo metterci la mano sul
fuoco, accadeva così spesso ed entrambe le parti erano così abituate,
che i bambini si sentivano incoraggiati, perché avevano imparato per
esperienza che i nostri soldati avevano l'ordine di non sparare loro. A
meno che non ci stessero sparando, cosa che finora non era mai successa,
che io sapessi.

La città era più degradata di quanto non lo fossero luoghi del genere un
tempo, credo che, prima del tentativo di riconquista del pianeta, i
criceti locali sperassero che la loro flotta avrebbe scacciato i
kristang e tutto sarebbe tornato alla normalità. Non fosse per una
grande forza di fastidiosi umani in residenza. Dopo il fallimento
dell'assalto, i criceti si erano resi conto che avevano perso il pianeta
per davvero, che sarebbero stati costretti ad andarsene e non sarebbero
mai più tornati, e avevano perso interesse a mantenere in ordine
qualcosa che si sarebbero lasciati alle spalle a uso e consumo dei
kristang. Piuttosto, erano risentiti, cosa che potevo capire, ed era
chiaro che progettavano di distruggere il posto un po' prima che
l'evacuazione fosse finita. Le principali infrastrutture, come il
Vettore e l'ascensore spaziale e i reattori a fusione, erano off limits,
secondo le strane regole di ingaggio che sembravano essere in vigore da
entrambe le parti. Dopotutto, se i ruhar avessero fatto saltare in aria
il Vettore sarebbe stato come ammettere che non avrebbero mai provato a
riconquistare Paradiso, e questo non avrebbe fatto bene al morale dei
criceti.

Stavamo attraversando l'area popolosa della città e ci stavamo
avvicinando al ponte occidentale, distante un paio di minuti, così
sembrava. Ero agitato per qualcosa che sentivo in procinto di accadere,
il linguaggio del corpo di Park e Olafson mi diceva che anche loro erano
pronti ad affrontare dei guai, ma era tutto tranquillo. Il ponte
occidentale era un'elegante struttura dall'aspetto delicato, molto più
lunga di quella del ponte orientale. Era stato costruito dai ruhar, ma
sui piloni di un precedente ponte kristang. Mi sentivo esposto mentre lo
attraversavamo, e tutti tirammo un sospiro di sollievo quando ci
ritrovammo all'asciutto dall'altra parte. Sporgendo la testa fuori dal
finestrino, rimasi a guardare il nostro convoglio finché l'ultimo camion
non ebbe superato il ponte. Stando alla mappa, e a quello che potevo
vedere, non c'era altro che terreni agricoli abbastanza pianeggianti,
che si estendevano per chilometri in tutte le direzioni.

«Planter, qui è Stinger.»

«Grazie per la scorta, Stinger, sei autorizzata a tornare alla base.»

«Ricevuto, Planter, divertitevi ora, coltivateci del cibo gustoso.
Stinger chiudo.» Fece volare il Pollo sopra di noi in un ampio cerchio,
poi ritrasse i pod delle armi, accelerò e salì verso est. Guardando il
Pollo volare, pensai che i nostri piloti, che su Paradiso guidavano
fighi aerei avanzati, dovevano amare quella missione. Mentre noialtri
volavamo su M4. Non sembrava giusto. Perché noi fanti non potevamo usare
i fucili sottratti ai ruhar, come i nostri aviatori guidavano i loro
aerei? Mi rimisi comodo al mio posto.

«Signore?» La voce di Olafson interruppe le mie fantasticherie. «Ha
incontrato i kristang, giusto? Abbiamo sentito che è andato alla loro
stazione quando è stato promosso.»

Una smorfia mi balenò in faccia prima che indossassi quella che speravo
fosse un'espressione neutra. Olafson se ne accorse, perché le sue
sopracciglia si alzarono. «C'è stata una cerimonia alla stazione, sì.
Io, ehm», mi arrabattai per coprire il mio errore, «il mio stomaco ha
fatto fatica con il cibo che hanno servito.» Era vero, anche se a farmi
star male era stato quello che aveva detto il kristang e il fatto di
avere dovuto ingoiare il mio orgoglio e tenere la bocca chiusa.

«Come sono, signore? Se posso chiederglielo. Non ne ho mai visto uno dal
vivo.» Doveva aver percepito il mio disagio di fronte a quella domanda:
«Dimentichi ciò che le ho chiesto, signore, quella roba è fuori dalla
mia giurisdizione\ldots».

Non lo avevo in alcun modo previsto. Un soldato nel penultimo camion
disse di avere visto una striscia appena prima dell'impatto: penso che
fosse la sua immaginazione. Potrebbe essere stato un cannone a rotaia,
oppure una specie di artiglieria silenziosa o un vero razzo senza fumo,
che non aveva lasciato un pennacchio di scarico. Qualunque cosa fosse,
uno di loro colpì i rimorchi di ogni camion, distruggendo il prezioso
fertilizzante e i semi, tutte cose insostituibili. La testata era
qualcosa di nuovo, poi generò più calore che esplosione: è probabile che
fosse un effetto calcolato per bruciare i nostri semi e i microrganismi
nel fertilizzante. I semi sparsi li avremmo potuti raccogliere, i semi
bruciati non sarebbero serviti più a nulla.

Erano proprio i semi che cercavano, per negare all'umanità la capacità
di sostentarsi su Paradiso, per costringere i kristang a spendere scarse
risorse per portare altri semi, o per riportare i loro cuccioli umani
sulla Terra. Nessuno, nessuno nei camion fu distrutto dal fuoco nemico,
il nostro unico morto fu un soldato sul sedile anteriore del Crivee che
viaggiava dietro l'ultimo camion; il loro veicolo tamponò il rimorchio
del mezzo che si era fermato all'improvviso e il conducente rimase
ferito quando la sua testa colpì il volante. Un rottame metallico
attraversò il cranio del soldato seduto sul sedile del passeggero,
uccidendolo sul colpo. Non credo che i criceti volessero ammazzare
qualcuno, il che era una magra consolazione per il morto e i feriti.

Uscimmo dai nostri veicoli in fretta per metterci al riparo, cercare i
nemici e soccorrere i feriti. Non avremmo mai scoperto che genere di
arma ci avesse colpito, perché Stinger reagì all'istante e inviò un paio
di missili ad alta velocità per colpire il punto da cui era stato
sferrato l'attacco. Quelle esplosioni ci indussero di nuovo a cercare
riparo, poi Stinger rombò sopra le nostre teste, con i pod delle armi
estesi e pronti ai guai.

«Planter, qui è Stinger, niente nemici in vista, passo.»

«Ricevuto, Stinger», risposi, mettendomi l'M4 a tracolla. «È probabile
che l'attacco fosse innescato da remoto.»

«E guidato», aggiunse Park, scrutando l'orizzonte nel mirino del fucile.
«Qualcuno aveva preso di mira i nostri rimorchi e nient'altro.»

Non aveva tutti i torti. O c'era un osservatore nascosto da qualche
parte, o le armi erano dotate di telecamere. Anche le armi aliene
intelligenti super-avanzate devono sapere cosa colpire. Stavo per dire
inutilmente a Stinger di mettersi alla ricerca di un osservatore, cosa
che stava già facendo, quando sul mio zPhone apparve un segnale
d'allarme. L'attacco al nostro convoglio non era stato l'unico
incidente, attacchi quasi simultanei erano stati lanciati in tutto il
continente. Chiamai il mio comandante e gli feci un rapido resoconto
della situazione, m'interruppe a metà del racconto visto che noi stavamo
bene per il momento, ma per altre unità non si poteva dire lo stesso.
Dal rumore in sottofondo, al quartier generale si era scatenato
l'inferno.

Un messaggio di testo dalla base aerea locale mi comunicò che era in
arrivo una Poiana per l'evacuazione medica, volevano che confermassi che
la zona di atterraggio era sicura. Merda, mi era sembrata sicura il
secondo prima che ci colpissero e sembrava sicura adesso. Il posto era
terreno agricolo pacifico, se ignoravi i pennacchi di fumo nero che
salivano dai nostri rimorchi bruciati e il Pollo in orbita attorno alla
zona con i pod delle armi caldi in cerca di guai e nella speranza di
trovarli. A parte questo, la zona era sicura. Risposi in modo
affermativo e ordinai ai miei di sparpagliarsi, abbastanza lontani l'uno
dall'altro da non costituire un singolo obiettivo, abbastanza vicini da
poterci sostenere a vicenda in una lotta. La Poiana arrivò e portò via
cinque feriti, noialtri facemmo inversione coi nostri Crivee, in modo da
poter tornare verso Habitrail. Stinger stava ancora volando in copertura
ad alta quota, questa volta non fece irruzione in città e disse che
poteva vedere dei criceti riunirsi nel campo da gioco di una scuola. Non
avevano armi in vista, e non ce n'erano più di una dozzina. Anche se i
criceti di Habitrail non erano coinvolti nell'attacco, dovevano sapere
che era successo qualcosa, i nostri rimorchi stavano ancora bruciando ed
emanando fumo nero che saliva alto nel cielo.

«Planter, qui è Crystal Fortress.» Crystal Fortress, ``Fortezza di
Cristallo'' era l'indicativo di chiamata del generale Maitland, il
comandante regionale. «Conferma di avere un morto in azione nella vostra
posizione, un chilometro a ovest di Habitrail?»

«Crystal Fortress, Planter conferma che abbiamo un morto in azione e
cinque feriti sono stati evacuati con una Poiana. Posizione a posto,
nessun segno di attività nemica qui, passo.» Speravo di avere altro
soccorso aereo, prima di tornare indietro attraversando la città.

«Planter, aspetta», ordinò la voce alla radio.

«Planter, ascolta bene.» La voce stavolta era quella del generale
Maitland in persona. «Abbiamo bisogno di mostrare ai criceti che questo
non è un gioco», la sua voce suonava strana, come se stesse leggendo
lentamente uno scritto, «che non possono colpirci e restare impuniti.
Andate a Habitrail e prendete otto ruhar a caso, a vostra discrezione.
Se qualcuno resiste, usate la forza letale. Allineate quegli otto ruhar
al centro della città e giustiziateli. Vogliamo che sia pubblico.»

Dovetti appoggiarmi al paraurti di un Crivee perché non mi cedessero le
gambe. La mia faccia doveva essere sbiancata, la gente mi guardava. Il
capitano Rivers mimò con le labbra in silenzio: «Che succede, capo?».
Rivolto a me.

\hfill\break

Ebbi un flashback di quando avevo otto anni ed ero a pesca in un
torrente con i miei amici, in una calda, umida giornata estiva, quel
genere di giornata estiva che capita così di rado nel Nord del Maine,
dove quando si arriva poco oltre i venti gradi i nativi si lamentano del
caldo. Io, i miei amici Bobby, Tommy e un bambino nuovo, Michael. Non
Mikey, Michael, ci aveva detto quando ci eravamo conosciuti. Era un
idiota, ma il padre era un pezzo grosso nella cartiera dove lavorava il
mio vecchio, quindi dovevamo essere gentili con lui e consentirgli di
venire con noi. Era del Wisconsin, e continuava a dirci quanto facesse
schifo il Maine, che eravamo stupidi a vivere lì, che vivevamo nel Maine
solo perché i nostri genitori non riuscivano a trovare lavoro altrove in
un posto decente.

Eravamo presso un torrente, a goderci quel genere di avventura boschiva,
da bambini allevati all'aperto senza supervisione, che tanto spaventa i
genitori di città, pescando e cercando salamandre. Michael aveva una
nuova canna da pesca con esche costose, noialtri avevamo attrezzatura
vecchia e usavamo pane arrotolato e pezzi di formaggio per esca. Bobby
catturò tre pesci di una certa dimensione, di quelli regolamentari per
la cena, e Tommy e io ci stavamo divertendo, ma senza catturare nulla
che valesse la pena di tenere. Michael era di cattivo umore perché non
aveva preso nulla, lamentandosi del fatto che i pesci nel torrente erano
stupidi, e noi stavamo alzando schizzi e gettando ombre che li
spaventavano. Per lo più lo ignoravamo, finché non catturò una rana per
sbaglio.

Torturò quella povera bestia. Già fu abbastanza brutto che ne avesse
agganciato la pancia e l'avesse squarciata tirandola a sé. Michael era
un piccolo, sadico pezzo di merda. Iniziò allegramente a usare un amo
per spellare la zampa della rana, mentre l'animale si dibatteva e lui lo
soffocava.

Non feci nulla per fermarlo. Me ne restai là in piedi, vergognandomi di
me stesso, con gli occhi fissi a terra. Anche il padre di Bobby lavorava
alla cartiera, ma Bobby era una persona migliore di me, pure a otto
anni. Lui afferrò la rana, la portò via e la calpestò su una roccia per
porre fine alle sue sofferenze. Michael ingaggiò una lotta cui Bobby e
Tommy misero fine molto in fretta, picchiandolo e spingendolo nel
torrente. Quando Michael urlò che l'avrebbe detto a suo padre, Bobby si
precipitò nel torrente e tenne la testa di quella piccola merda
nell'acqua fangosa finché non piagnucolò: «Mi arrendo». Prima che ce ne
andassimo, Tommy si ruppe la nuova canna da pesca di Michael sulle
ginocchia e gettò la sua costosa cassetta delle esche nelle acque
profonde del torrente. E Bobby avvisò Michael che se l'avesse detto a
suo padre, o se l'avessimo rivisto fuori da scuola, l'avremmo pestato a
sangue. A otto anni, quella minaccia significava che facevamo sul serio,
e Michael lo sapeva.

Non rivedemmo Michael mai più. Suo padre si trasferì in un mulino in
Oregon quel settembre e portò la famiglia con sé. Michael, m'immaginavo,
crescendo sarebbe diventato uno che picchia la moglie o un serial
killer, o entrambi, due facce della stessa medaglia.

La mia mente deve aver avuto un flashback di quel momento perché allora
me ne ero rimasto a guardare e non avevo fatto nulla per difendere un
innocente. Non sarebbe mai più accaduto. E a maggior ragione non con
indosso un'uniforme dell'esercito.

\hfill\break

«Crystal Fortress», scandii molto lentamente, a voce abbastanza alta
perché i soldati attorno a me non potessero perdersi quello che stavo
dicendo, «conferma che mi stai ordinando di uccidere otto civili.» La
mascella inferiore mi tremava mentre parlavo: «Giustiziare otto civili
in pubblico». Avevo messo il mio zPhone in vivavoce e Maitland sapeva
cosa stavo facendo.

«Cazzo, colonnello Bishop. Non mi piace più di quanto piaccia a te. Ce
la stiamo cavando a stento qui; questa merda non può andare avanti.
Questi sono gli ordini. Eseguili.»

«Signore», dissi fissando un soldato dopo l'altro negli occhi, «non
posso farlo.» Gli sguardi che ricevetti erano scioccati quanto me. Anche
se avessi dato l'ordine, era improbabile che qualcuno del convoglio
avrebbe obbedito. In che casino ci eravamo cacciati? Merda, non avremmo
mai dovuto lasciare la Terra. Gli umani non avevano alcun motivo di
stare lì.

«Maledizione, Bishop, non indossi quell'uccello sull'uniforme per
bellezza.»

«Signore, con tutto il rispetto, lei non ha l'autorità per scavalcare le
regole di ingaggio, o il codice di condotta.» Perché cazzo stavo avendo
quella conversazione? «Possiamo\ldots»

Una strana voce s'intromise, una voce che riconobbi essere quella di un
kristang che usava un traduttore. «Colonnello Joseph Bishop, ti stai
rifiutando di obbedire a un ordine diretto?» Nemmeno il traduttore
riuscì a rimuovere tutto il sibilo lucertolesco dalla voce.

Lucertole. Alzai gli occhi senza volerlo. Avevano navi in orbita che
avrebbero potuto vaporizzare la mia intera squadra lì. Dovevo fare
attenzione con le vite dei venti umani sotto il mio comando, molta
attenzione. Spensi il mio zPhone per un momento, infilandolo in una
tasca. In fretta, spiegai la situazione ai soldati attorno a me: «Le
lucertole mi hanno mostrato chi sono davvero, quando ero lassù per la
cerimonia di promozione. Le voci che avete sentito rispetto ai messaggi
contenuti nei biscotti della fortuna sono vere, le lucertole non hanno
salvato la Terra dai ruhar, ma hanno cacciato i ruhar in modo da poter
violentare il nostro pianeta indisturbate. Siamo fottuti in ogni caso,
l'unica cosa che ci è rimasta è la nostra umanità, e non ho intenzione
di rinunciarci. Le lucertole mi ordinano di uccidere otto criceti a
caso».

«Perché otto, signore?», chiese Olafson. Pensai non fosse esattamente la
domanda fondamentale da porre in quel momento.

Tenendo in alto le mani, presi a illustrare: «Le lucertole hanno quattro
dita per mano, pensano in termini di otto, non dieci come noi. Dovremmo
uccidere otto civili per ogni umano ucciso. Quando i nazisti occuparono
Roma e gli italiani opposero resistenza, i primi spararono a dieci o
venti civili per ogni tedesco ucciso. Presero dei civili a caso per la
strada, li misero in fila su un ponte, li fucilarono e gettarono i loro
corpi nel fiume. L'esercito degli Stati Uniti non lo fa». Sull'argomento
avevo un po' di confusione rispetto ai dati di fatto, ma gli elementi di
base erano corretti. Con lo zPhone di nuovo acceso, risposi: «In quanto
soldato dell'esercito degli Stati Uniti, sono tenuto a rifiutare ordini
illegali. Assassinare a sangue freddo dei civili è illegale».

Ci fu uno strillo kristang inferocito che non si traduceva, poi il mio
zPhone si spense. Intendo del tutto, niente luci, le lucertole dovevano
averlo disattivato.

«E adesso, signore?», chiese il capitano Rivers.

Come diavolo facevo a saperlo? Niente nel mio addestramento mi aveva
preparato a quel casino di situazione. «Stacchiamo il camion di testa da
quel rimorchio rotto.» La parte posteriore della cabina del camion era
bruciata e le finestre erano saltate, ma le solide gomme erano ancora
rotonde e funzionali. Gli altri due furgoni erano in cattive condizioni,
le loro celle a combustibile avevano preso fuoco per il calore. «Guiderò
il camion da solo, se le lucertole di sopra decidono di farmi fuori per
disobbedienza, non voglio portare nessun altro con me. E non discutete,
è un ordine. Partiamo, torniamo alla base, teniamo i Crivee
sparpagliati.» Anche se le lucertole potevano colpirci dall'orbita a
prescindere da quanto ampiamente fossimo sparpagliati, non volevo
facilitare loro le cose. Forse avremmo dovuto mettere tutti gli zPhone
in un sacchetto che avrei tenuto con me, così le lucertole non avrebbero
potuto prenderci di mira sulla base del segnale\ldots{}

Rivers mise una mano sul suo auricolare, poi me lo porse assieme al suo
zPhone: «Stinger vuole parlare con lei, signore». Così non avevano
disattivato tutti i nostri zPhone, soltanto il mio.

«Stinger, qui è Planter, vai pure.»

«Ho ascoltato le sue comunicazioni, signore. I kristang mi hanno
ordinato di prendere di mira una scuola di Habitrail. Mi sono rifiutata
di obbedire all'ordine.» La disobbedienza stava diventando popolare.

«Ricevuto, Stinger, grazie. Stiamo per partire\ldots»

«Merda!», gridò Stinger. «Ho perso il controllo! I comandi non
rispondono! Il mitragliere non riesce a bloccare le armi!»

«Sparpagliatevi!», sbraitai, e indicai i Polli in avvicinamento a bassa
quota da sud-ovest, con i pod delle armi estesi. «Le lucertole hanno il
controllo della cannoniera!»

Tutti correvano; mi tuffai in un fosso sul ciglio della strada, ma mi
resi conto in un lampo di quanto fosse stupido e inutile. Ed egoista. Le
lucertole avevano un problema con me, non con la mia squadra. Perciò mi
alzai e agitai le braccia verso il Pollo, cercando di attirare
l'attenzione. Il mio piano, se così si può chiamare, era di sollevare il
dito medio dopo che le lucertole mi avessero lanciato contro un missile.
Sarebbe stato stupido anche quello, è probabile che le lucertole non
riconoscessero gesti umani volgari.

Non fosse che non attaccarono me. I due missili restanti del Pollo
schizzarono fuori dalle rotaie, sopra le nostre teste, e colpirono la
scuola con una fiammeggiante esplosione. «Non siamo stati noi!», gridò
Stinger in tono febbrile. «Oh merda! Hanno interrotto la corrente!
Stiamo precipitando!»

Mentre lo guardavo, il Pollo barcollò in aria, poi cadde come una pietra
da circa novanta metri di altitudine, fracassandosi a terra con grande
violenza e tracciando un solco nel terreno a circa quattrocento metri a
sud dal punto in cui mi trovavo. Senza che dicessi nulla, i soldati
iniziarono a correre sul luogo dell'incidente, io arrivai per primo
perché non ero appesantito da tutto l'armamentario di battaglia come gli
altri.

Il Pollo era molto meno danneggiato di quanto mi aspettassi, con
qualsiasi cosa i criceti lo avessero costruito era roba resistente. La
coda e le due alette si erano spezzate, il tutto si era cappottato di
trecentosessanta gradi, ma la cabina di pilotaggio e il quadro elettrico
dietro i sedili erano per lo più interi. A differenza dei due occupanti
umani. Il mitragliere sul sedile anteriore era senza testa, una lama
della ventola del motore si era staccata e l'aveva tagliata di netto. Al
pilota dietro di lui, di cui lessi il nome sulla tuta di volo -- si
chiamava Collins -- a lei mancava la parte inferiore del braccio
sinistro e le usciva sangue dalla bocca. Non era cosciente, con
attenzione le slacciai la cintura e sollevai il casco, che una grande
crepa tagliava da cima a fondo. Fu terribile.

Maledizione. Come diavolo faceva a esserci una ragazza, voglio dire una
donna, così giovane alla guida di un Pollo? Per quanto malconcia fosse,
tra i due sembravo comunque io il più vecchio.

«Uh», gemette e aprì un occhio.

«Ssst. Non ti muovere, Collins. Sono il colonnello Bishop, sono
Planter.»

«Non riuscivo a fermarlo, signore, non riuscivo a fermarlo», la sua voce
si smorzò, sebbene le sue labbra si muovessero ancora in silenzio,
mentre cercava di parlarmi.

«Va tutto bene, resta con me. Collins? Collins?» La testa le si
afflosciò in avanti mentre la tenevo tra le braccia, un'ultima bolla di
sangue le uscì dalla bocca e morì. Le sue ultime parole non avevano
espresso paura, mi aveva detto di avere fatto il suo dovere, che non
aveva ucciso criceti nella scuola. In qualche modo, per via di una
combinazione di circostanze, Collins era passata dall'essere una
bambina, all'arruolarsi nell'esercito, alla scuola di volo in
elicottero, al volontariato per il servizio con l'Unef, qualificandosi
per pilotare una cannoniera Pollo, fino a qui. A morire tra le braccia
di un perfetto sconosciuto. Uccisa da una specie aliena che non aveva un
suo concetto di umanità. Non c'era una buona ragione perché la sua vita
finisse lì, quel giorno. Crollai e piansi.

Non ero il solo a stare in piedi con le spalle abbassate, singhiozzando
piano. Il tenente Collins non era stata l'unica persona che avevamo
perso quel giorno, la sua morte era stata la goccia che aveva fatto
traboccare il vaso. Non era stata uccisa in combattimento, era stata
uccisa da creature che avrebbero dovuto essere nostri alleati, nostri
protettori, i nostri salvatori. L'intera situazione su Paradiso era
andata a puttane così in fretta che avevamo perso quasi ogni speranza
nel giro di pochi mesi.

Rivers mi toccò una spalla, mi guardò negli occhi e scosse la testa una
volta, in silenzio. Messaggio ricevuto. Il superiore doveva dare
l'esempio. Mi raddrizzai, asciugandomi con rabbia gli occhi con la parte
posteriore di una manica. Riconsegnando lo zPhone a Rivers, gli dissi di
fare lui un resoconto della situazione al quartier generale dell'Unef,
perché non pensavo di riuscire a essere professionale in quel momento.

Nessuno di noi sapeva cos'altro dire. Non c'era niente da dire. Dopo un
po', tirammo fuori i due corpi dal Pollo e li caricammo sul sedile
posteriore del camion che guidavo io. Dopo avere attraversato il ponte
occidentale, invece di seguire la strada che tagliava Habitrail, diedi
ordine di attraversare i campi in un ampio cerchio, cosa che ora
potevamo fare, non essendo più appesantiti dai rimorchi. Ci volle un'ora
per aggirare la città, il che comportava anche dare la caccia ai
percorsi giusti per guadare torrenti, sobbalzare attraverso terra e
fango; il lavoro e la concentrazione mantenevano la gente assorbita dal
compito di tornare alla base. Io restai indietro col camion, mantenendo
dall'ultimo Crivee una distanza tale che, se le lucertole avessero
deciso di farmi fuori con un colpo di cannone a rotaia o un missile, non
ci sarebbero stati danni collaterali tra la mia gente. Aveva anche senso
che fossero i Crivee ad andare in ricognizione del percorso migliore per
attraversare il paese, perché il camion non era altrettanto funzionale
fuori strada.

Eravamo tornati sul percorso principale da meno di mezz'ora, quando la
colonna di testa sostò e io mi fermai quattrocento metri dietro di loro.
Rivers mi chiamò sullo zPhone che avevo preso in prestito: «Il comando
dice di aspettare qui, signore, stanno mandando una Poiana». Il comando
non aveva detto il motivo. Una sola Poiana non sarebbe bastata per
evacuare tutti noi, e non potevo credere che l'Unef volesse abbandonare
i Crivee funzionanti. Forse la situazione che ci aspettava era così
brutta che il comando ci stava portando rinforzi con armi più pesanti?
Provai a chiamare il mio zPhone, ma poteva connettersi solo a Rivers.
Ecco cosa succedeva quando un'altra specie controllava tutte le tue
comunicazioni. Non avremmo dovuto lasciare che accadesse.

Naturalmente, un'altra specie controllava anche tutto il nostro cibo e
le altre provviste, e la posizione strategica e la Terra, quindi le
comunicazioni su Paradiso avrebbero potuto essere l'ultimo dei nostri
problemi.

Una Poiana, scortata da un Pollo, arrivò rombando, ci volteggiò sopra e
atterrò sulla strada di fronte alla colonna. Vidi Rivers andare avanti a
piedi per parlare con i soldati. Il mio zPhone suonò, era Rivers:
«Signore, la vogliono qui». C'era una sfumatura di allerta nella sua
voce. Lasciai il mio M4 nel camion e mi avviai di buon passo. Mentre
passavo accanto ai Crivee, i soldati mi facevano il saluto, alcuni
avevano le lacrime agli occhi, sembravano tutti arrabbiati. Che cazzo
stava succedendo? Raggiunsi il gruppo riunito attorno alla Poiana e
salutai un certo capitano Randolph. La Poiana era dell'esercito degli
Stati Uniti, come tutte le truppe che trasportava.

«Capitano Randolph.»

Mi fece il saluto militare, mi fissò un istante negli occhi, ma non
riuscì a non distogliere lo sguardo. Qualunque cosa stesse succedendo,
mi metteva a disagio. «Colonnello Bishop, ho l'ordine di dichiararla in
arresto. Consegni la sua arma da fianco, prego.» Tese la mano, quasi in
segno di scusa.

«Arresto?» Ero davvero scioccato. «Con quale accusa?», chiesi.

«Rifiuto di obbedire agli ordini di un superiore, signore.»

«Rifiuto di eseguire ordini illegali, capitano.»

«Signore, questo è al di sopra della mia giurisdizione», disse l'altro a
fatica. «Colonnello, lei non è l'unica persona di cui il quartier
generale dell'Unef ha ordinato l'arresto, si dice che un sacco di unità
abbiano rifiutato quegli ordini. Abbiamo sentito», abbassò la voce, «che
i kristang hanno colpito alcune unità dall'orbita quando hanno rifiutato
gli ordini direttamente da, ehm\ldots», indicò il cielo con il pollice.
«La prego, signore, non posso rischiare la vita dei miei uomini mancando
di arrestarla.» Questo implicava che le lucertole stavano guardando e ci
avrebbero ucciso tutti se non mi fossi arreso.

«Colonnello», iniziò a dire Rivers, lanciando uno sguardo significativo
al dito che teneva sulla sicura del suo M4.

«Capitano Rivers, sei tu al comando qui. Riporta queste persone alla
base in sicurezza, evita problemi lungo la strada», ordinai. Tenendo la
mia arma da fianco con due dita, la consegnai a Randolph. «Mi portate al
quartier generale dell'Unef?»

«No, signore», il suo pomo d'Adamo fece su e giù, nervosamente, «tutti i
prigionieri sono diretti a una base delle luc\ldots{} dei kristang,
signore.»